\input{configTD}

\title{Analyse Numérique TD2 --- Systèmes linéaires, LU et conditionnement}

\author{}
\date{}

\begin{document}
\sffamily
\maketitle

% ====================================================================
%  RAPPELS DE COURS : SYSTÈMES LINÉAIRES, GAUSS, LU
% ====================================================================

\begin{mdframed}[backgroundcolor=ThemeLight, linecolor=Theme,
  linewidth=1.5pt, innertopmargin=10pt, innerbottommargin=10pt]
\textbf{\large Rappels --- Résolution de systèmes linéaires}

\medskip

\textbf{Problème.}
On cherche $x \in \mathbb{R}^n$ tel que $Ax = b$,
où $A$ est une matrice $n \times n$ inversible
et $b \in \mathbb{R}^n$.

\medskip

\textbf{Systèmes triangulaires.}
Si $A$ est \emph{triangulaire supérieure}, on résout
par \textbf{substitution arrière} :
\[
  x_n = \frac{b_n}{a_{nn}}\,,
  \qquad
  x_i = \frac{1}{a_{ii}}
    \Bigl(b_i - \sum_{j=i+1}^{n} a_{ij}\,x_j\Bigr),
  \quad i = n\!-\!1, \ldots, 1.
\]
Si $A$ est \emph{triangulaire inférieure}, on résout
par \textbf{substitution avant} (de $i=1$ vers $i=n$).

\medskip

\textbf{Pivot de Gauss.}
On transforme le système $Ax=b$ en un système
triangulaire supérieur $Ux=c$ par des opérations
élémentaires sur les lignes :
$L_i \leftarrow L_i - m_{ik}\,L_k$
($m_{ik}$ est appelé \emph{multiplicateur}).

\medskip

\textbf{Factorisation $A = LU$.}
Les multiplicateurs de Gauss forment une matrice
$L$ triangulaire inférieure avec des $1$ sur la diagonale.
On obtient
\[
  A = L\,U
\]
où $U$ est la matrice triangulaire supérieure obtenue.
La résolution se fait alors en \textbf{deux étapes} :
\[
  L\,y = b \quad \text{(substitution avant)}
  \qquad \text{puis} \qquad
  U\,x = y \quad \text{(substitution arrière)}.
\]

\textbf{Avantage.}
Si le second membre $b$ change mais pas $A$,
on réutilise $L$ et $U$ sans refaire l'élimination.
\end{mdframed}

% ====================================================================
\section*{Exercice 1 --- Construction guidée : pivot de Gauss et LU}
% ====================================================================

On considère le système linéaire $Ax = b$ avec
\[
  A = \begin{pmatrix}
    4 & -1 & 0\\
    -1 & 4 & -1\\
    0 & -1 & 4
  \end{pmatrix},
  \qquad
  b = \begin{pmatrix} 7\\ 6\\ 7 \end{pmatrix}.
\]
Cette matrice tridiagonale apparaît comme matrice de rigidité
simplifiée lors de la discrétisation d'une poutre.

\begin{enumerate}

% ---------- Étape 1 : élimination ----------
\item \textbf{Première étape de Gauss.}
Pour éliminer le coefficient $-1$ en position $(2,1)$ :
\begin{enumerate}
  \item Quel multiplicateur $m_{21}$ faut-il utiliser
        dans l'opération $L_2 \leftarrow L_2 - m_{21}\,L_1$ ?
  \item Effectuer cette opération sur la matrice augmentée
        $(A \mid b)$.
\end{enumerate}

\ifthenelse{\boolean{showSolutions}}{
  \begin{solution}
    $m_{21} = \dfrac{a_{21}}{a_{11}} = \dfrac{-1}{4}$.
    L'opération $L_2 \leftarrow L_2 - (-\tfrac{1}{4})\,L_1
    = L_2 + \tfrac{1}{4}\,L_1$ donne
    \[
      \left(\begin{array}{ccc|c}
        4 & -1 & 0 & 7\\[4pt]
        0 & \tfrac{15}{4} & -1 & \tfrac{31}{4}\\[4pt]
        0 & -1 & 4 & 7
      \end{array}\right).
    \]
  \end{solution}
}{}

\item \textbf{Deuxième étape de Gauss.}
Déterminer le multiplicateur $m_{32}$ et effectuer
l'opération $L_3 \leftarrow L_3 - m_{32}\,L_2$
pour obtenir un système triangulaire supérieur.

\ifthenelse{\boolean{showSolutions}}{
  \begin{solution}
    $m_{32} = \dfrac{-1}{15/4} = -\dfrac{4}{15}$.
    L'opération $L_3 \leftarrow L_3 + \tfrac{4}{15}\,L_2$ donne
    \[
      \left(\begin{array}{ccc|c}
        4 & -1 & 0 & 7\\[4pt]
        0 & \tfrac{15}{4} & -1 & \tfrac{31}{4}\\[4pt]
        0 & 0 & \tfrac{56}{15} & \tfrac{139}{15}
      \end{array}\right).
    \]
  \end{solution}
}{}

% ---------- Résolution du système triangulaire ----------
\item Résoudre le système triangulaire supérieur
obtenu par substitution arrière.

\ifthenelse{\boolean{showSolutions}}{
  \begin{solution}
    \[
      x_3 = \frac{139/15}{56/15} = \frac{139}{56}\,,
      \qquad
      x_2 = \frac{4}{15}\Bigl(\frac{31}{4} + x_3\Bigr)
           = \frac{19}{7}\,,
      \qquad
      x_1 = \frac{1}{4}(7 + x_2) = \frac{17}{7}\,.
    \]
    Donc
    $x = \bigl(\tfrac{17}{7},\;\tfrac{19}{7},\;
    \tfrac{139}{56}\bigr)^T$.
  \end{solution}
}{}

% ---------- Extraction de L et U ----------
\item \textbf{Factorisation LU.}
Les multiplicateurs trouvés aux questions 1 et 2
sont $m_{21} = -\tfrac{1}{4}$ et $m_{32} = -\tfrac{4}{15}$.
Former les matrices $L$ et $U$ telles que $A = LU$.

\ifthenelse{\boolean{showSolutions}}{
  \begin{solution}
    \[
      L = \begin{pmatrix}
        1 & 0 & 0\\[3pt]
        -\tfrac{1}{4} & 1 & 0\\[3pt]
        0 & -\tfrac{4}{15} & 1
      \end{pmatrix},
      \qquad
      U = \begin{pmatrix}
        4 & -1 & 0\\[3pt]
        0 & \tfrac{15}{4} & -1\\[3pt]
        0 & 0 & \tfrac{56}{15}
      \end{pmatrix}.
    \]
    On peut vérifier que $LU = A$.
  \end{solution}
}{}

% ---------- Résolution via LU ----------
\item Retrouver la solution $x$ en résolvant
successivement $Ly = b$ puis $Ux = y$.

\ifthenelse{\boolean{showSolutions}}{
  \begin{solution}
    Substitution avant ($Ly = b$) :
    \[
      y_1 = 7,\qquad
      y_2 = 6 + \tfrac{1}{4} \cdot 7 = \tfrac{31}{4},\qquad
      y_3 = 7 + \tfrac{4}{15} \cdot \tfrac{31}{4} = \tfrac{139}{15}.
    \]
    Substitution arrière ($Ux = y$) :
    on retrouve
    $x = \bigl(\tfrac{17}{7},\;\tfrac{19}{7},\;
    \tfrac{139}{56}\bigr)^T$.
  \end{solution}
}{}

\end{enumerate}

% ====================================================================
\section*{Exercice 2 --- Effet d'une perturbation du second membre}
% ====================================================================

On reprend la matrice $A$ et sa factorisation $LU$
de l'exercice précédent.
On remplace le second membre par
\[
  \tilde{b} = \begin{pmatrix} 7\\ 6{,}1\\ 7 \end{pmatrix}.
\]

\begin{enumerate}
\item Sans refaire l'élimination de Gauss,
expliquer comment obtenir rapidement la nouvelle solution
$\tilde{x}$.

\ifthenelse{\boolean{showSolutions}}{
  \begin{solution}
    La factorisation $A = LU$ ne dépend pas de $b$.
    Il suffit de résoudre
    $L\tilde{y} = \tilde{b}$ puis $U\tilde{x} = \tilde{y}$.
  \end{solution}
}{}

\item Résoudre le système $A\tilde{x} = \tilde{b}$.

\ifthenelse{\boolean{showSolutions}}{
  \begin{solution}
    Substitution avant :
    \[
      \tilde{y}_1 = 7,\qquad
      \tilde{y}_2 = 6{,}1 + \tfrac{7}{4} = \tfrac{157}{20},\qquad
      \tilde{y}_3 = 7 + \tfrac{4}{15}\times \tfrac{157}{20}
                   = \tfrac{683}{75}.
    \]
    Substitution arrière :
    \[
      \tilde{x}_3 = \tfrac{683}{280} \approx 2{,}439,\qquad
      \tilde{x}_2 \approx 2{,}839,\qquad
      \tilde{x}_1 \approx 2{,}460.
    \]
  \end{solution}
}{}

\item Comparer $x$ et $\tilde{x}$.
La perturbation de $0{,}1$ sur $b_2$ produit-elle
une grande variation de la solution ?

\ifthenelse{\boolean{showSolutions}}{
  \begin{solution}
    On avait $x \approx (2{,}429;\;2{,}714;\;2{,}482)^T$.
    Les composantes varient de l'ordre de $0{,}03$ à $0{,}13$ :
    la variation reste modérée.
    Le système n'est pas particulièrement sensible
    aux perturbations sur cet exemple.
  \end{solution}
}{}
\end{enumerate}

% ====================================================================
%  RAPPELS DE COURS : CONDITIONNEMENT
% ====================================================================

\begin{mdframed}[backgroundcolor=ThemeLight, linecolor=Theme,
  linewidth=1.5pt, innertopmargin=10pt, innerbottommargin=10pt]
\textbf{\large Rappels --- Conditionnement d'un système linéaire}

\medskip

\textbf{Norme infinie d'une matrice.}
\[
  \|A\|_{\infty}
  = \max_{1 \leq i \leq n}
    \sum_{j=1}^{n} |a_{ij}|
  \qquad
  \text{(plus grande somme des valeurs absolues d'une ligne).}
\]

\textbf{Nombre de condition.}
\[
  \kappa(A)
  = \|A\|\;\|A^{-1}\|.
\]

\textbf{Interprétation.}
Si $Ax = b$ et $A\tilde{x} = \tilde{b}$
avec $\tilde{b} = b + \delta b$, alors
\[
  \frac{\|\tilde{x} - x\|}{\|x\|}
  \;\leq\;
  \kappa(A)\;\frac{\|\delta b\|}{\|b\|}\,.
\]
Un $\kappa$ grand signifie qu'une \emph{petite} erreur
relative sur $b$ peut produire une \emph{grande} erreur
relative sur $x$ : le système est \textbf{mal conditionné}.

\medskip

\textbf{Ordres de grandeur.}
$\kappa \approx 1$ : bien conditionné ;
$\kappa \approx 10^k$ : on peut perdre
environ $k$ chiffres significatifs.
\end{mdframed}

% ====================================================================
\section*{Exercice 3 --- Système mal conditionné}
% ====================================================================

On considère la matrice
\[
  B = \begin{pmatrix}
    1 & 1\\
    1 & 1{,}01
  \end{pmatrix}.
\]

\begin{enumerate}
\item Calculer $\det(B)$ puis $B^{-1}$.

\ifthenelse{\boolean{showSolutions}}{
  \begin{solution}
    $\det(B) = 1 \times 1{,}01 - 1 \times 1 = 0{,}01$.
    \[
      B^{-1} = \frac{1}{0{,}01}
      \begin{pmatrix} 1{,}01 & -1\\ -1 & 1 \end{pmatrix}
      = \begin{pmatrix} 101 & -100\\ -100 & 100 \end{pmatrix}.
    \]
  \end{solution}
}{}

\item Calculer $\|B\|_{\infty}$ et $\|B^{-1}\|_{\infty}$.

\ifthenelse{\boolean{showSolutions}}{
  \begin{solution}
    $\|B\|_{\infty} = \max(1+1,\;1+1{,}01) = 2{,}01$.
    \quad
    $\|B^{-1}\|_{\infty} = \max(101+100,\;100+100) = 201$.
  \end{solution}
}{}

\item En déduire $\kappa_{\infty}(B)$ et interpréter.

\ifthenelse{\boolean{showSolutions}}{
  \begin{solution}
    \[
      \kappa_{\infty}(B) = 2{,}01 \times 201 \approx 404.
    \]
    Le système est mal conditionné :
    une erreur relative de $1\,\%$ sur $b$
    peut produire une erreur relative
    de l'ordre de $4\,\%$ sur $x$ (facteur $\approx 400$).
  \end{solution}
}{}

\item On résout $Bx = (2,\;2{,}01)^T$
puis $B\tilde{x} = (2,\;2{,}02)^T$.
Calculer les deux solutions et vérifier
que l'amplification de l'erreur est conforme
à la borne donnée par $\kappa$.

\ifthenelse{\boolean{showSolutions}}{
  \begin{solution}
    $x = B^{-1}(2,\;2{,}01)^T
    = (101 \times 2 - 100 \times 2{,}01,\;
       -100 \times 2 + 100 \times 2{,}01)^T
    = (1,\;1)^T$.

    $\tilde{x} = B^{-1}(2,\;2{,}02)^T
    = (101 \times 2 - 100 \times 2{,}02,\;
       -100 \times 2 + 100 \times 2{,}02)^T
    = (0,\;2)^T$.

    \medskip

    La perturbation relative sur $b$ :
    $\|\delta b\|_\infty / \|b\|_\infty
    = 0{,}01 / 2{,}02 \approx 0{,}5\,\%$.

    La variation relative sur $x$ :
    $\|\tilde{x}-x\|_\infty / \|x\|_\infty
    = 1 / 1 = 100\,\%$.

    Le facteur d'amplification est $\approx 200$,
    inférieur à $\kappa \approx 404$ : la borne est respectée.
  \end{solution}
}{}
\end{enumerate}

% ====================================================================
\section*{Exercice 4 --- Application à un modèle discrétisé}
% ====================================================================

On cherche à approcher la solution du problème
\[
  -u''(x) = 1,
  \qquad x \in [0,1],
  \qquad u(0) = u(1) = 0
\]
par différences finies avec un pas $h = 1/4$.
On place les n\oe uds intérieurs
$x_1 = \tfrac{1}{4}$, $x_2 = \tfrac{1}{2}$,
$x_3 = \tfrac{3}{4}$,
et on note $u_i \approx u(x_i)$.

\begin{enumerate}
\item Écrire l'approximation centrée de $u''(x_i)$
en fonction de $u_{i-1}$, $u_i$ et $u_{i+1}$.

\ifthenelse{\boolean{showSolutions}}{
  \begin{solution}
    \[
      u''(x_i) \approx
      \frac{u_{i+1} - 2u_i + u_{i-1}}{h^2}\,.
    \]
  \end{solution}
}{}

\item En remplaçant dans $-u'' = 1$
et en utilisant $u_0 = u_4 = 0$,
écrire le système linéaire vérifié par $u_1, u_2, u_3$.

\ifthenelse{\boolean{showSolutions}}{
  \begin{solution}
    Avec $h^2 = 1/16$, l'équation
    $-\dfrac{u_{i+1}-2u_i+u_{i-1}}{h^2}=1$
    devient $-u_{i+1}+2u_i-u_{i-1} = \tfrac{1}{16}$.
    Le système s'écrit
    \[
      \underbrace{\begin{pmatrix}
        2 & -1 & 0\\
        -1 & 2 & -1\\
        0 & -1 & 2
      \end{pmatrix}}_{A_h}
      \begin{pmatrix} u_1\\ u_2\\ u_3 \end{pmatrix}
      =
      \begin{pmatrix} 1/16\\ 1/16\\ 1/16 \end{pmatrix}.
    \]
  \end{solution}
}{}

\item Résoudre ce système (par Gauss ou LU).

\ifthenelse{\boolean{showSolutions}}{
  \begin{solution}
    Par Gauss (ou en exploitant la symétrie) :
    \[
      u_1 = \tfrac{3}{32},
      \qquad
      u_2 = \tfrac{1}{8},
      \qquad
      u_3 = \tfrac{3}{32}.
    \]
  \end{solution}
}{}

\item La solution exacte est
$u(x) = \dfrac{x(1-x)}{2}$.
Comparer les valeurs numériques aux n\oe uds.

\ifthenelse{\boolean{showSolutions}}{
  \begin{solution}
    \[
      u\bigl(\tfrac{1}{4}\bigr) = \tfrac{3}{32},
      \quad
      u\bigl(\tfrac{1}{2}\bigr) = \tfrac{1}{8},
      \quad
      u\bigl(\tfrac{3}{4}\bigr) = \tfrac{3}{32}.
    \]
    Les valeurs numériques coïncident exactement
    avec la solution.
    Ce résultat remarquable s'explique
    par le fait que $u$ est un polynôme de degré $2$
    et que la différence finie centrée est exacte
    pour les polynômes de degré $\leq 3$.
  \end{solution}
}{}
\end{enumerate}

\end{document}
